\subsection{Right Triangles} 
If we have a triangle with corners at $A$, $B$, and $C$, we write $\triangle ABC$. 
\begin{icpic}
\draw[geom] (-2.5,0) -- (-1.5,2) -- (2.5,0) -- cycle;
\filldraw (-2.5,0) circle (0.1) node[anchor=north]{$A$};
\filldraw (-1.5,2) circle (0.1) node[anchor=south]{$B$};
\filldraw (2.5,0) circle (0.1) node[anchor=north]{$C$};
\draw (-1.68,1.64) -- ++(0.36,-0.18) -- ++(0.18,0.36);
\end{icpic}
In this diagram:
\begin{itemize}
\item We see a small square at $\angle B$ (``\makeblank'').
\item This marks a \makeblank[1in] angle.
\item So $\triangle ABC$ is a \makeblank[1in] triangle.
\item $\overline{AC}$ is called the \makeblank[1.5in] of the triangle.
\item $\overline{AB}$ and $\overline{BC}$ are called the \makeblank[1in] of the triangle.
\end{itemize}
If we have a diagonal line segment in a coordinate plane (one that isn't \makeblank[1in] or \makeblank[1in]) there is a natural way to draw a right triangle with that segment as the hypotenuse.
\begin{Example}
Plot $A=(-2,3)$ and $B=(3,-1)$. Choose a point $C$ so that $\triangle ABC$ is a right triangle with hypotenuse $\overline{AB}$.
\begin{icpic}
\setgraphsquare{5}
\GraphLarge
\end{icpic}
\begin{itemize}
\item We could put $C$ at \blankpair or \blankpair
\item If $A=(x_A,y_A)$ and $B=(x_B,y_B)$, we could put $C$ at \blankpair or \blankpair.
\end{itemize}
\end{Example}
